\documentclass[12pt,a4paper]{article}

\usepackage[T2A]{fontenc}
\usepackage[utf8]{inputenc}
\usepackage[russian]{babel}
\usepackage{amsmath,amssymb,amsfonts}
\usepackage{geometry}
\usepackage{booktabs}
\usepackage{enumitem}
\usepackage{array}

\geometry{margin=2.2cm}

\title{Финальные математические версии ИИС: V4, V5, V6}
\author{Проект анализа интегрального индекса состояния}
\date{\today}

\begin{document}
\maketitle

\section{Назначение документа}
Настоящий документ фиксирует три актуальные негормональные математические версии модели ИИС:
\begin{itemize}[leftmargin=1.25cm]
    \item \textbf{IISVersion4} --- базовая калиброванная физиологическая версия без гормонального блока;
    \item \textbf{IISVersion5} --- контрастная нелинейная версия, уменьшающая плоскость шкалы;
    \item \textbf{IISVersion6} --- универсальная gated-версия, моделирующая состояние как мягкую смесь нескольких режимов.
\end{itemize}

Во всех трёх версиях используются только реально измеряемые или допустимые для proxy-режима блоки:
\[
A,\ \Gamma,\ V,\ Q.
\]
Гормональные блоки \(H\) и \(K\) в этих версиях исключены.

\section{Общие обозначения}
Используются следующие обозначения:
\begin{itemize}[leftmargin=1.25cm]
    \item \(P_L, P_R\) --- суммарная мощность левого и правого EEG-блоков;
    \item \(\alpha_{\mathrm{asym}}\) --- alpha-asymmetry;
    \item \(P_\gamma\) --- gamma power;
    \item \(\gamma/\alpha\) --- gamma/alpha ratio;
    \item \(HR\) --- heart rate;
    \item \(HF/LF\) --- отношение высокочастотной и низкочастотной компонент HRV;
    \item \(\sigma(x)=\dfrac{1}{1+e^{-x}}\) --- логистическая сигмоида;
    \item \(\Gamma_+=\max(\Gamma,0)\);
    \item \(V_+=\max(V,0)\);
    \item \(\varepsilon>0\) --- малый стабилизирующий параметр.
\end{itemize}

Используется мягкая центрированная нормировка
\[
T(x;c,w)=\tanh\!\left(\frac{x-c}{w}\right),
\]
а также знако-сохраняющая степенная нелинейность
\[
\phi_p(x)=\operatorname{sign}(x)\,|x|^p.
\]

Для многорежимной версии используется softmax:
\[
\operatorname{softmax}(z_1,\dots,z_n)_i=
\frac{e^{z_i}}{\sum_{j=1}^{n}e^{z_j}}.
\]

\section{Общие базовые компоненты}
\subsection{Компонент \(A\)}
Сначала alpha-асимметрия:
\[
A_{\alpha}=T(\alpha_{\mathrm{asym}};\ 0.000613,\ 0.000530).
\]

Затем суммарная межполушарная асимметрия:
\[
A_{\mathrm{total}}=
T\!\left(
\frac{P_L-P_R}{P_L+P_R+\varepsilon};\
-0.0299,\ 0.0287
\right).
\]

Итоговый компонент:
\[
A=\tanh\!\left(1.10\,A_{\alpha}+0.40\,A_{\mathrm{total}}\right).
\]

\subsection{Компонент \(\Gamma\)}
Сначала gamma-мощность переводится в лог-шкалу:
\[
G_{\log}=\ln(1+P_\gamma\cdot 10^{12}).
\]

Далее:
\[
G_{\gamma}=T(G_{\log};\ 1.543,\ 0.850),
\]
\[
G_{\gamma/\alpha}=T(\ln(1+\gamma/\alpha);\ 0.000367,\ 0.000434).
\]

Итоговый gamma-компонент:
\[
\Gamma=\tanh\!\left(0.30\,G_{\gamma}+0.70\,G_{\gamma/\alpha}\right).
\]

\subsection{Компонент \(V\)}
Компонент по ЧСС:
\[
V_{hr}=T(118.740-HR;\ 0,\ 46.718).
\]

Компонент по HRV:
\[
V_{hrv}=T(\ln(1+HF/LF);\ 0.968,\ 1.222).
\]

Итоговый вегетативный компонент:
\[
V=\tanh\!\left(0.80\,V_{hr}+0.20\,V_{hrv}\right).
\]

\section{Версия IISVersion4}
\subsection{Компонент \(Q_4\)}
В версии 4 интегральное качество состояния считается как простая физиологическая интеграция:
\[
Q_4=\tanh\!\left(0.55A+0.45V-0.25\Gamma\right).
\]

В proxy-режиме используется мягкая якоризация на self-report:
\[
Q_{4,\mathrm{anchor}}=
\tanh\!\left(
2.2\left(
0.62\,valence_{\mathrm{norm}}
0.38(1-arousal_{\mathrm{norm}})
-0.5
\right)\right),
\]
\[
Q_{4,\mathrm{proxy}}=0.65\,Q_4+0.35\,Q_{4,\mathrm{anchor}}.
\]

\subsection{Итоговый индекс \(IIS_4\)}
Статическая формула версии 4:
\[
IIS_4=\sigma\!\left(0.10A-0.05\Gamma+0.25V+0.60Q_4\right).
\]

\section{Версия IISVersion5}
\subsection{Компонент \(Q_5\)}
В версии 5 компонент \(Q\) становится более нелинейным и учитывает не только сумму блоков, но и их согласованность.

Линейное ядро:
\[
Q_{\mathrm{lin}}=0.50A+0.42V-0.30\Gamma.
\]

Когерентность:
\[
Q_{\mathrm{coh}}=A\cdot V-0.50|A-V|-0.18\Gamma_+.
\]

Нелинейно преобразованные переменные:
\[
\widehat{A}=\phi_{0.82}(A),\qquad
\widehat{V}=\phi_{0.82}(V),\qquad
\widehat{\Gamma}=\phi_{1.05}(\Gamma).
\]

Синергия:
\[
S_{AV}=\operatorname{sign}(\widehat{A}+\widehat{V})\sqrt{|\widehat{A}\widehat{V}|+\varepsilon}.
\]

Энергия:
\[
E_{AV}=\frac{|\widehat{A}|+|\widehat{V}|}{2}.
\]

Конфликт:
\[
C_{AV}=|\widehat{A}-\widehat{V}|+0.16\,\widehat{\Gamma}_+\,(1-\widehat{V}_+),
\qquad
\widehat{\Gamma}_+=\max(\widehat{\Gamma},0),\quad
\widehat{V}_+=\max(\widehat{V},0).
\]

Итоговая формула:
\[
Q_5=
\tanh\!\Bigl(
1.4\,Q_{\mathrm{lin}}
0.55\,Q_{\mathrm{coh}}
0.42\,S_{AV}
0.28\,E_{AV}\operatorname{sign}(Q_{\mathrm{lin}})
-0.24\,C_{AV}
\Bigr).
\]

В proxy-режиме:
\[
Q_{5,\mathrm{anchor}}=
\tanh\!\left(
2.4\left(
0.65\,valence_{\mathrm{norm}}
0.35(1-arousal_{\mathrm{norm}})
-0.5
\right)\right),
\]
\[
Q_{5,\mathrm{proxy}}=0.75Q_5+0.25Q_{5,\mathrm{anchor}}.
\]

\subsection{Итоговый индекс \(IIS_5\)}
Базовое линейное ядро:
\[
core_5=0.12A-0.07\Gamma+0.24V+0.57Q_5.
\]

Нелинейно преобразованные переменные:
\[
\widetilde{A}=\phi_{0.84}(A),\qquad
\widetilde{V}=\phi_{0.84}(V),\qquad
\widetilde{Q}=\phi_{0.84}(Q_5).
\]

Центрированное ядро:
\[
z_5=0.14\widetilde{A}-0.08\Gamma+0.22\widetilde{V}+0.56\widetilde{Q}-0.08.
\]

Синергия режимов:
\[
S_{\mathrm{reg}}=
\operatorname{sign}(\widetilde{A}+\widetilde{V})\sqrt{|\widetilde{A}\widetilde{V}|+\varepsilon}
+0.8\,\operatorname{sign}(\widetilde{Q}+\widetilde{V})\sqrt{|\widetilde{Q}\widetilde{V}|+\varepsilon}.
\]

Энергия фазы:
\[
E_{\mathrm{phase}}=\sqrt{\frac{\widetilde{A}^2+\widetilde{V}^2+\widetilde{Q}^2}{3}}.
\]

Конфликт:
\[
C_{\mathrm{phase}}=|\widetilde{A}-\widetilde{V}|+0.75|\widetilde{Q}-\widetilde{V}|.
\]

Контрастный слой:
\[
\begin{aligned}
contrast_5={}&
2.6z_5
+0.8z_5^3
+0.18(Q_5-V)
-0.06\Gamma_+ \\
&+0.52S_{\mathrm{reg}}
+0.34E_{\mathrm{phase}}\operatorname{sign}(z_5)
-0.22C_{\mathrm{phase}}
-0.10\Gamma_+(1-V_+).
\end{aligned}
\]

Базовая и контрастная части:
\[
IIS_{5,\mathrm{base}}=\sigma\!\left(core_5+0.22S_{\mathrm{reg}}-0.10C_{\mathrm{phase}}\right),
\]
\[
IIS_{5,\mathrm{contrast}}=0.5+0.5\tanh(contrast_5).
\]

Финальная смесь:
\[
IIS_5=\operatorname{clip}\!\left(
0.65\,IIS_{5,\mathrm{base}}+0.35\,IIS_{5,\mathrm{contrast}},
0,1\right).
\]

\section{Версия IISVersion6}
\subsection{Идея версии 6}
Версия 6 интерпретирует состояние не как результат одной глобальной формулы, а как мягкую смесь трёх режимов:
\begin{itemize}[leftmargin=1.25cm]
    \item \textbf{regulation} --- согласованное регулируемое состояние;
    \item \textbf{mobilization} --- напряжённая мобилизация;
    \item \textbf{depletion} --- истощающий режим.
\end{itemize}

\subsection{Нелинейно преобразованные компоненты}
\[
\widehat{A}=\phi_{0.78}(A),\qquad
\widehat{V}=\phi_{0.78}(V),\qquad
\widehat{\Gamma}=\phi_{1.05}(\Gamma),\qquad
\widehat{Q}=\phi_{0.78}(Q_6),
\]
\[
\widehat{\Gamma}_+=\max(\widehat{\Gamma},0).
\]

\subsection{Гейты для \(Q_6\)}
Сначала вычисляются логиты режимов:
\[
z^{(Q)}_{\mathrm{reg}}=0.95\widehat{A}+1.05\widehat{V}-0.55\widehat{\Gamma}_+,
\]
\[
z^{(Q)}_{\mathrm{mob}}=-0.18\widehat{A}+0.48\widehat{V}-1.05\widehat{\Gamma}_+,
\]
\[
z^{(Q)}_{\mathrm{dep}}=-0.80\widehat{A}-0.95\widehat{V}+0.72\widehat{\Gamma}_+.
\]

Гейты:
\[
\bigl(g^{(Q)}_{\mathrm{reg}},g^{(Q)}_{\mathrm{mob}},g^{(Q)}_{\mathrm{dep}}\bigr)
=
\operatorname{softmax}\!\left(
1.65\,z^{(Q)}_{\mathrm{reg}},
1.65\,z^{(Q)}_{\mathrm{mob}},
1.65\,z^{(Q)}_{\mathrm{dep}}
\right).
\]

\subsection{Локальные уровни режима для \(Q_6\)}
\[
S_{AV}^{(6)}=\operatorname{sign}(\widehat{A}+\widehat{V})\sqrt{|\widehat{A}\widehat{V}|+\varepsilon}.
\]

\[
q_{\mathrm{reg}}=
0.76+0.18\tanh\!\left(
1.10\widehat{A}+1.15\widehat{V}-0.45\widehat{\Gamma}_++0.30S_{AV}^{(6)}
\right),
\]
\[
q_{\mathrm{mob}}=
0.50+0.14\tanh\!\left(
0.10\widehat{A}+0.55\widehat{V}-1.20\widehat{\Gamma}_+
\right),
\]
\[
q_{\mathrm{dep}}=
0.18+0.16\tanh\!\left(
-0.65|\widehat{A}|-0.95|\widehat{V}|+0.95\widehat{\Gamma}_+
\right).
\]

Итоговый \(Q_6\):
\[
Q_6=
\operatorname{clip}\!\left(
g^{(Q)}_{\mathrm{reg}}q_{\mathrm{reg}}
+g^{(Q)}_{\mathrm{mob}}q_{\mathrm{mob}}
+g^{(Q)}_{\mathrm{dep}}q_{\mathrm{dep}},
0,1
\right).
\]

В proxy-режиме:
\[
Q_{6,\mathrm{anchor}}=0.70\,valence_{\mathrm{norm}}+0.30(1-arousal_{\mathrm{norm}}),
\]
\[
Q_{6,\mathrm{proxy}}=0.82Q_6+0.18Q_{6,\mathrm{anchor}}.
\]

\subsection{Гейты итогового индекса \(IIS_6\)}
Логиты итоговых режимов:
\[
z_{\mathrm{reg}}=0.60\widehat{Q}+0.45\widehat{V}+0.15\widehat{A}-0.30\widehat{\Gamma}_+,
\]
\[
z_{\mathrm{mob}}=-0.05\widehat{Q}+0.20\widehat{V}-0.55\widehat{\Gamma}_+-0.08\widehat{A},
\]
\[
z_{\mathrm{dep}}=-0.70\widehat{Q}-0.55\widehat{V}-0.20\widehat{A}+0.35\widehat{\Gamma}_+.
\]

\[
\bigl(g_{\mathrm{reg}},g_{\mathrm{mob}},g_{\mathrm{dep}}\bigr)=
\operatorname{softmax}\!\left(
1.65\,z_{\mathrm{reg}},
1.65\,z_{\mathrm{mob}},
1.65\,z_{\mathrm{dep}}
\right).
\]

\subsection{Локальные уровни режима}
\[
S_{QV}^{(6)}=\operatorname{sign}(\widehat{Q}+\widehat{V})\sqrt{|\widehat{Q}\widehat{V}|+\varepsilon},
\]
\[
C_{QV}^{(6)}=|\widehat{A}-\widehat{V}|+0.75|\widehat{Q}-\widehat{V}|.
\]

\[
i_{\mathrm{reg}}=
0.80+0.16\tanh\!\left(
0.75\widehat{Q}+0.55\widehat{V}+0.18\widehat{A}-0.25\widehat{\Gamma}_++0.25S_{QV}^{(6)}
\right),
\]
\[
i_{\mathrm{mob}}=
0.54+0.12\tanh\!\left(
-0.08\widehat{Q}+0.32\widehat{V}-0.52\widehat{\Gamma}_++0.08\widehat{A}
\right),
\]
\[
i_{\mathrm{dep}}=
0.20+0.16\tanh\!\left(
-0.62\widehat{Q}-0.52\widehat{V}-0.18\widehat{A}+0.38\widehat{\Gamma}_+
\right).
\]

Промежуточная смесь:
\[
R_6=
g_{\mathrm{reg}}i_{\mathrm{reg}}
+g_{\mathrm{mob}}i_{\mathrm{mob}}
+g_{\mathrm{dep}}i_{\mathrm{dep}}.
\]

\subsection{Энтропия переходности и финальная формула}
Нормированная энтропия гейта:
\[
H(g)=
-\frac{g_{\mathrm{reg}}\ln g_{\mathrm{reg}}
+g_{\mathrm{mob}}\ln g_{\mathrm{mob}}
+g_{\mathrm{dep}}\ln g_{\mathrm{dep}}}{\ln 3}.
\]

Баланс между regulation и depletion:
\[
B_6=g_{\mathrm{reg}}-g_{\mathrm{dep}}.
\]

Итоговый индекс:
\[
IIS_6=
\operatorname{clip}\!\left(
R_6
-0.08\,H(g)
+0.05\,B_6
-0.10\,C_{QV}^{(6)},
0,1
\right).
\]

Для служебной диагностики также используется вспомогательное ядро:
\[
raw_6=0.14\widehat{A}-0.08\widehat{\Gamma}+0.24\widehat{V}+0.54\widehat{Q}.
\]

\section{Краткая интерпретация трёх версий}
\begin{center}
\begin{tabular}{>{\raggedright\arraybackslash}p{2.6cm} >{\raggedright\arraybackslash}p{4.2cm} >{\raggedright\arraybackslash}p{7.2cm}}
\toprule
Версия & Основная идея & Главная особенность \\
\midrule
V4 & Базовая негормональная физиологическая версия & Почти минимальная и хорошо интерпретируемая формула, удобная как опорная модель \\
V5 & Контрастная нелинейная версия & Усиливает середину и края шкалы, лучше раскрывает геометрию без смены архитектуры \\
V6 & Универсальная gated-версия & Моделирует состояние как мягкую смесь режимов regulation, mobilization и depletion \\
\bottomrule
\end{tabular}
\end{center}

\section{Итог}
Формально версии \(V4\), \(V5\) и \(V6\) образуют последовательность:
\[
\text{простая физиологическая интеграция}
\;\longrightarrow\;
\text{контрастная нелинейная интеграция}
\;\longrightarrow\;
\text{мягкая многорежимная gated-модель}.
\]

Таким образом:
\begin{itemize}[leftmargin=1.25cm]
    \item \textbf{V4} --- базовый калиброванный эталон;
    \item \textbf{V5} --- шаг к более богатой нелинейной геометрии;
    \item \textbf{V6} --- наиболее универсальная тестовая архитектура среди текущих негормональных версий.
\end{itemize}

\end{document}
